\documentclass[11pt,a4paper]{report}
\usepackage[utf8]{inputenc}
\usepackage[T1]{fontenc}
\usepackage[french]{babel}
\usepackage[a4paper, margin=2.2cm, top=2.5cm, bottom=2.5cm]{geometry}
\usepackage{xcolor}
\usepackage{tcolorbox}
\usepackage{booktabs}
\usepackage{tabularx}
\usepackage{fancyhdr}
\usepackage{lastpage}
\usepackage{enumitem}
\usepackage{amsmath, amssymb}
\usepackage{titlesec}
\usepackage{hyperref}
\usepackage{listings}
\usepackage{tikz}
\usetikzlibrary{shapes.geometric, arrows.meta, positioning, calc, backgrounds, fit}

\definecolor{primaryNavy}{RGB}{26, 54, 93}
\definecolor{accentTeal}{RGB}{13, 148, 136}
\definecolor{accentGold}{RGB}{217, 119, 6}
\definecolor{darkSlate}{RGB}{30, 41, 59}
\definecolor{lightBg}{RGB}{248, 250, 252}
\definecolor{borderColor}{RGB}{226, 232, 240}

\titleformat{\chapter}[display]
  {\normalfont\huge\bfseries\color{primaryNavy}}{\chaptertitlename\ \thechapter}{10pt}{\Huge}
\titleformat{\section}
  {\normalfont\Large\bfseries\color{primaryNavy}}{\thesection}{1em}{}
\titleformat{\subsection}
  {\normalfont\large\bfseries\color{accentTeal}}{\thesubsection}{1em}{}
\titleformat{\subsubsection}
  {\normalfont\normalsize\bfseries\color{darkSlate}}{\thesubsubsection}{1em}{}

\pagestyle{fancy}
\fancyhf{}
\fancyhead[L]{\small\textcolor{darkSlate}{\textbf{MatOS} --- Architecture Système \& Module Biométrique IA}}
\fancyhead[R]{\small\textcolor{accentTeal}{\nouppercase{\leftmark}}}
\fancyfoot[L]{\small\textcolor{gray}{Dossier de Spécification Technique \& Algorithmique --- ISO/IEC 30107-3}}
\fancyfoot[R]{\small\textcolor{darkSlate}{\textbf{Page \thepage\ sur \pageref{LastPage}}}}
\renewcommand{\headrulewidth}{0.6pt}
\renewcommand{\footrulewidth}{0.4pt}

\tcbset{
  basebox/.style={
    colback=lightBg,
    colframe=borderColor,
    boxrule=0.8pt,
    arc=4pt,
    left=10pt, right=10pt, top=8pt, bottom=8pt,
    fonttitle=\bfseries\color{white}
  },
  execbox/.style={
    basebox,
    colback=primaryNavy!5,
    colframe=primaryNavy,
    title=Synthèse d'Ingénierie Système
  },
  alertbox/.style={
    basebox,
    colback=accentGold!8,
    colframe=accentGold,
    title=Contrainte Critique de Sécurité \& Conformité
  },
  techbox/.style={
    basebox,
    colback=accentTeal!6,
    colframe=accentTeal,
    title=Spécification Algorithmique \& Standard
  }
}

\hypersetup{
    colorlinks=true,
    linkcolor=primaryNavy,
    filecolor=accentTeal,      
    urlcolor=accentTeal,
    citecolor=primaryNavy,
}

\lstset{
  backgroundcolor=\color{lightBg},
  basicstyle=\ttfamily\footnotesize\color{darkSlate},
  breaklines=true,
  captionpos=b,
  commentstyle=\color{gray},
  keywordstyle=\bfseries\color{primaryNavy},
  stringstyle=\color{accentTeal},
  frame=single,
  rulecolor=\color{borderColor}
}

\begin{document}

\begin{titlepage}
    \centering
    \vspace*{1.5cm}
    {\Huge\bfseries\color{primaryNavy} MatOS}\\[0.4cm]
    {\LARGE\bfseries\color{accentTeal} Marketplace Universelle de Location de Matériel Sécurisée}\\[1.5cm]
    
    \rule{\linewidth}{1.8pt}\\[0.4cm]
    {\huge\bfseries ARCHITECTURE TECHNIQUE GLOBALE, MODULE BIOMÉTRIQUE IA\\\& SÉCURISATION MONÉTIQUE}\\[0.3cm]
    {\large Spécifications ISO/IEC 30107-3 (PAD), Algorithmes Anti-Deepfake, Chaîne de Preuve RFC 3161 \& Schéma SQL ACID}\\[0.4cm]
    \rule{\linewidth}{1.8pt}\\[2.2cm]
    
    \begin{minipage}{0.45\textwidth}
        \begin{flushleft} \large
            \textbf{Auteur :}\\
            Direction Technique (CTO)\\
            Pôle IA \& Sécurité des Systèmes
        \end{flushleft}
    \end{minipage}
    \hfill
    \begin{minipage}{0.45\textwidth}
        \begin{flushright} \large
            \textbf{Statut :}\\
            Spécification Technique Officielle\\
            \textbf{Date :} Août 2026
        \end{flushright}
    \end{minipage}
    
    \vfill
    {\small Normes de Référence : ISO/IEC 30107-3, NIST SP 800-63B, RFC 3161 (TSP), CNDP (Loi 09-08)}
\end{titlepage}

\tableofcontents
\newpage

\begin{tcolorbox}[execbox]
Ce document présente les spécifications d'ingénierie logicielle, d'infrastructure cloud et de vision par ordinateur pour la plateforme \textbf{MatOS}. L'architecture est conçue pour résoudre le défi fondamental de la confiance dans la location de matériel à haute valeur ajoutée sans dépendre de tiers de vérification coûteux ni compromettre la fluidité de l'expérience utilisateur mobile.

Le système articule une stack technologique moderne combinant un front-end hybride (\textbf{Next.js 14} pour l'indexation SEO publique et \textbf{Flutter} pour l'accès caméra natif), un back-end asynchrone ultra-rapide sous \textbf{FastAPI} orchestrant une base de données \textbf{PostgreSQL 16 / PostGIS}, et un moteur d'intelligence artificielle propriétaire exécutant en mémoire vive le pipeline de contrôle d'identité conforme à la norme \textbf{ISO/IEC 30107-3} (OCR de la CIN marocaine, test de vivacité 3D actif, analyse spectrale anti-deepfake et comparaison faciale 512-D par ArcFace).
\end{tcolorbox}

\chapter{Architecture Système \& Infrastructure Distribuée}

\section{Topologie Globale \& Découpage Modulaire}
La plateforme MatOS s'articule autour d'une architecture orientée services à couplage faible, garantissant une résilience opérationnelle totale et une extensibilité horizontale sans interruption de service.

\begin{figure}[h!]
\centering
\begin{tikzpicture}[node distance=1.3cm and 1.8cm, auto, >=Latex, every node/.style={font=\scriptsize}]
    % Client Layer
    \node[draw=primaryNavy, fill=primaryNavy!10, rounded corners, minimum width=3.2cm, minimum height=1.3cm, align=center] (clientWeb) {\textbf{Vitrine Web (Next.js 14)}\\SSR / ISR / SEO Google\\PWA Ultra-Légère};
    \node[draw=primaryNavy, fill=primaryNavy!10, rounded corners, minimum width=3.2cm, minimum height=1.3cm, align=center, below=0.6cm of clientWeb] (clientMobile) {\textbf{Application Mobile (Flutter)}\\iOS \& Android Natif\\Flux Caméra 60 fps};
    
    % Gateway
    \node[draw=accentTeal, fill=accentTeal!15, thick, rounded corners, minimum width=3.4cm, minimum height=3cm, right=1.6cm of clientWeb, yshift=-0.9cm, align=center] (gateway) {\textbf{API Gateway (FastAPI)}\\Auth JWT RS256 Asymétrique\\Rate Limiting Redis 7\\Validation Pydantic v2\\Orchestrateur Asynchrone};

    % Services
    \node[draw=darkSlate, fill=lightBg, rounded corners, minimum width=3.6cm, minimum height=1.1cm, right=1.6cm of gateway, yshift=1.5cm, align=center] (iaService) {\textbf{Moteur IA Biométrique}\\MediaPipe Face Mesh + ONNX\\Liveness PAD ISO 30107-3};
    \node[draw=darkSlate, fill=lightBg, rounded corners, minimum width=3.6cm, minimum height=1.1cm, right=1.6cm of gateway, yshift=0cm, align=center] (bizService) {\textbf{Services Métier}\\Contrats DOC + Escrow CMI\\Messagerie Chiffrée Socket};
    \node[draw=darkSlate, fill=lightBg, rounded corners, minimum width=3.6cm, minimum height=1.1cm, right=1.6cm of gateway, yshift=-1.5cm, align=center] (storageService) {\textbf{Horodatage \& Médias}\\Hash SHA-256 + RFC 3161\\Cloudflare R2 Object Storage};

    % Data Layer
    \node[draw=accentGold, fill=accentGold!15, thick, rounded corners, minimum width=4.5cm, minimum height=1.3cm, below=1.2cm of gateway, align=center] (db) {\textbf{Base de Données Principale}\\PostgreSQL 16 + PostGIS (Geo) + pgcrypto\\Chiffrement AES-256 des données CNDP};

    % Connections
    \draw[->, thick, primaryNavy] (clientWeb.east) -- (gateway.west |- clientWeb.east);
    \draw[->, thick, primaryNavy] (clientMobile.east) -- (gateway.west |- clientMobile.east);
    
    \draw[->, thick, accentTeal] (gateway.east) -- (iaService.west);
    \draw[->, thick, accentTeal] (gateway.east) -- (bizService.west);
    \draw[->, thick, accentTeal] (gateway.east) -- (storageService.west);
    
    \draw[<->, thick, accentGold] (gateway.south) -- (db.north);
\end{tikzpicture}
\caption{Schéma de l'architecture distribuée de la plateforme MatOS}
\end{figure}

\section{Justification des Choix Technologiques}
\begin{enumerate}[leftmargin=*]
    \item \textbf{Front-End Découplé (Next.js 14 \& Flutter) :} Le référencement naturel (SEO) étant le vecteur d'acquisition principal à coût marginal nul, Next.js 14 génère des pages HTML statiques réhydratées pour chaque annonce d'équipement (ex. \textit{« Location perforateur Bosch Casablanca Maârif »}). L'application mobile Flutter assure quant à elle le contrôle granulaire de la caméra pour interdire l'injection de fichiers falsifiés et fluidifier les états des lieux guidés.
    \item \textbf{Back-End FastAPI \& Python Asynchrone :} FastAPI offre des performances brutes comparables à Node.js et Go grâce à son moteur \textit{uvloop}, tout en permettant l'intégration directe et sans surcoût de sérialisation des modèles de deep learning Python (PyTorch, ONNX Runtime, OpenCV).
    \item \textbf{PostgreSQL 16 \& PostGIS :} La recherche de matériel repose sur une indexation spatiale géodésique (\textit{R-Tree indexing}) permettant d'interroger les objets disponibles dans un rayon de 5 à 20 kilomètres autour de la position GPS de l'utilisateur avec un temps d'exécution inférieur à 5 millisecondes.
\end{enumerate}

\chapter{Pipeline Biométrique IA \& Détection du Vivant (ISO/IEC 30107-3)}

\section{Algorithmes de Liveness Check Actif \& Détection 3D}
Pour contrer les attaques par présentation (masques en silicone, photographies imprimées, vidéos rejouées sur écran haute définition), le module de vivacité impose un défi interactif aléatoire :

\begin{figure}[h!]
\centering
\begin{tikzpicture}[node distance=0.8cm and 1cm, auto, >=Latex, every node/.style={font=\scriptsize}]
    % Pipeline Nodes
    \node[draw=primaryNavy, fill=lightBg, rounded corners, minimum width=2.6cm, minimum height=1.3cm, align=center] (step1) {\textbf{1. Capture CIN}\\Extraction OCR\\Contrôle MRZ Modulo};
    \node[draw=accentTeal, fill=lightBg, rounded corners, minimum width=2.9cm, minimum height=1.3cm, right=0.7cm of step1, align=center] (step2) {\textbf{2. Liveness Check}\\Caméra Live Native\\Challenge 3D (EAR $<0.2$)};
    \node[draw=accentGold, fill=lightBg, rounded corners, minimum width=2.9cm, minimum height=1.3cm, right=0.7cm of step2, align=center] (step3) {\textbf{3. Anti-Deepfake}\\Analyse Spectrale FFT\\Reflets Cornéens};
    \node[draw=primaryNavy, fill=accentTeal!20, thick, rounded corners, minimum width=2.6cm, minimum height=1.3cm, right=0.7cm of step3, align=center] (step4) {\textbf{4. Matching}\\ArcFace 512-D\\$\text{Score} \ge 0,78$};

    % Arrows
    \draw[->, thick, primaryNavy] (step1) -- (step2);
    \draw[->, thick, accentTeal] (step2) -- (step3);
    \draw[->, thick, accentGold] (step3) -- (step4);
\end{tikzpicture}
\caption{Pipeline séquentiel de validation biométrique et détection d'attaque (PAD)}
\end{figure}

Le calcul de l'ouverture palpébrale pour la validation du clignement naturel est régi par le ratio d'aspect de l'œil (\textbf{Eye Aspect Ratio - EAR}) calculé sur les repères 3D extraits par MediaPipe (repères $p_1$ à $p_6$) :
\begin{equation}
\text{EAR} = \frac{\|p_2 - p_6\|_2 + \|p_3 - p_5\|_2}{2 \|p_1 - p_4\|_2}
\end{equation}
Un clignement biologique valide produit une chute transitoire de la valeur avec $\text{EAR} < 0,20$ pendant une durée caractéristique comprise entre 100 et 350 millisecondes.

\section{Module Spectral Anti-Deepfake \& Extraction d'Embeddings ArcFace}
Les avatars synthétiques et les flux vidéo générés en temps réel par des modèles de diffusion ou GANs sont neutralisés par une analyse fréquentielle bidimensionnelle par Transformée de Fourier Rapide (\textbf{FFT}) :
\begin{equation}
F(u, v) = \sum_{x=0}^{M-1} \sum_{y=0}^{N-1} f(x, y) \cdot e^{-j 2\pi \left(\frac{ux}{M} + \frac{vy}{N}\right)}
\end{equation}
L'absence de micro-textures de peau naturelles et la présence de périodicités anormales dans les hautes fréquences spatiales entraînent le rejet automatique du flux.

La comparaison entre le portrait extrait de la CIN marocaine et le visage validé par le liveness check est opérée par un réseau de neurones convolutif profond \textbf{ArcFace} (marge angulaire additive). Les deux vecteurs d'embedding normalisés de dimension 512 ($\vec{u}, \vec{v} \in \mathbb{R}^{512}$) sont comparés par similarité cosinus :
\begin{equation}
\text{Score}(\vec{u}, \vec{v}) = \cos(\theta) = \frac{\vec{u} \cdot \vec{v}}{\|\vec{u}\|_2 \cdot \|\vec{v}\|_2}
\end{equation}
Le seuil de décision est fixé à $\mathbf{0,78}$, garantissant un Taux de Fausse Acceptation (\textit{False Acceptance Rate - FAR}) strictement inférieur à $0,001\%$ et un Taux de Faux Rejet (\textit{False Rejection Rate - FRR}) inférieur à $1,2\%$.

\chapter{Chaîne de Preuve Numérique \& Horodatage RFC 3161}

\section{Protocole Cryptographique d'État des Lieux Contradictoire}
L'état du matériel lors de la prise en charge et lors de la restitution fait l'objet d'un scellement cryptographique irréfragable garantissant son opposabilité juridique complète devant les tribunaux :

\begin{figure}[h!]
\centering
\begin{tikzpicture}[node distance=1.2cm and 1.5cm, auto, >=Latex, every node/.style={font=\scriptsize}]
    \node[draw=primaryNavy, fill=lightBg, rounded corners, minimum width=3cm, minimum height=1.2cm, align=center] (capture) {\textbf{1. Clichés Guidés}\\4 à 8 photos AR\\Exifs GPS + Appareil};
    \node[draw=accentTeal, fill=lightBg, rounded corners, minimum width=3.2cm, minimum height=1.2cm, right=1cm of capture, align=center] (hash) {\textbf{2. Hachage SHA-256}\\Génération Hash Unique\\JSON Canonique Signé};
    \node[draw=accentGold, fill=lightBg, rounded corners, minimum width=3.2cm, minimum height=1.2cm, right=1cm of hash, align=center] (tsp) {\textbf{3. Horodatage RFC 3161}\\Autorité Qualifiée\\Jeton Infalsifiable};

    \draw[->, thick, primaryNavy] (capture) -- (hash);
    \draw[->, thick, accentTeal] (hash) -- (tsp);
\end{tikzpicture}
\caption{Workflow de génération de la chaîne de preuve numérique pour l'état des lieux}
\end{figure}

Le fichier d'état des lieux contient la concaténation canonique des empreintes numériques :
\begin{equation}
H_{\text{global}} = \text{SHA-256}\left(\text{ID}_{\text{contrat}} \,\|\, \text{ID}_{\text{bailleur}} \,\|\, \text{ID}_{\text{locataire}} \,\|\, \text{GPS} \,\|\, \sum_{i=1}^n \text{SHA-256}(\text{Photo}_i)\right)
\end{equation}
Ce condensat est scellé par l'autorité d'horodatage qualifiée, rendant toute modification ultérieure techniquement impossible et immédiatement détectable.

\chapter{Schéma Relationnel de Données (PostgreSQL 16 DDL)}

\section{Définition des Tables Principales \& Contraintes d'Intégrité}
Le modèle de données assure la cohérence transactionnelle ACID et le respect des règles d'anonymisation de la CNDP :

\begin{lstlisting}[language=SQL, caption={DDL des tables maîtresses de la plateforme MatOS}]
CREATE EXTENSION IF NOT EXISTS "uuid-ossp";
CREATE EXTENSION IF NOT EXISTS "postgis";
CREATE EXTENSION IF NOT EXISTS "pgcrypto";

CREATE TABLE users (
    id UUID PRIMARY KEY DEFAULT uuid_generate_v4(),
    phone_number VARCHAR(20) UNIQUE NOT NULL,
    cin_number VARCHAR(20) UNIQUE,
    full_name VARCHAR(100),
    is_cin_verified BOOLEAN DEFAULT FALSE,
    biometric_embedding BYTEA, -- Chiffre AES-256 via pgcrypto
    subscription_tier VARCHAR(20) DEFAULT 'FREE', -- FREE, PREMIUM, PRO
    created_at TIMESTAMP WITH TIME ZONE DEFAULT NOW()
);

CREATE TABLE items (
    id UUID PRIMARY KEY DEFAULT uuid_generate_v4(),
    owner_id UUID NOT NULL REFERENCES users(id) ON DELETE CASCADE,
    title VARCHAR(150) NOT NULL,
    description TEXT,
    category VARCHAR(50) NOT NULL,
    replacement_value NUMERIC(10,2) NOT NULL,
    daily_price NUMERIC(10,2) NOT NULL,
    location GEOGRAPHY(Point, 4326) NOT NULL,
    is_available BOOLEAN DEFAULT TRUE,
    created_at TIMESTAMP WITH TIME ZONE DEFAULT NOW()
);

CREATE TABLE rentals (
    id UUID PRIMARY KEY DEFAULT uuid_generate_v4(),
    item_id UUID NOT NULL REFERENCES items(id),
    renter_id UUID NOT NULL REFERENCES users(id),
    start_date DATE NOT NULL,
    end_date DATE NOT NULL,
    total_amount NUMERIC(10,2) NOT NULL,
    commission_amount NUMERIC(10,2) NOT NULL,
    cmi_auth_token VARCHAR(100),
    status VARCHAR(30) DEFAULT 'PENDING_KYC',
    handover_hash_rfc3161 VARCHAR(64),
    return_hash_rfc3161 VARCHAR(64),
    created_at TIMESTAMP WITH TIME ZONE DEFAULT NOW()
);

CREATE INDEX idx_items_location ON items USING GIST(location);
CREATE INDEX idx_rentals_status ON rentals(status);
\end{lstlisting}

\chapter*{Bibliographie \& Références Techniques}
\addcontentsline{toc}{chapter}{Bibliographie \& Références Techniques}

\begin{enumerate}[leftmargin=*]
    \item \textbf{ISO/IEC :} \textit{ISO/IEC 30107-3:2017 --- Information Technology — Biometric Presentation Attack Detection}.
    \item \textbf{Deng, J., Guo, J., Xue, N., \& Zafeiriou, S. :} \textit{ArcFace: Additive Angular Margin Loss for Deep Face Recognition}, IEEE CVPR, 2019.
    \item \textbf{Internet Engineering Task Force (IETF) :} \textit{RFC 3161 --- Internet X.509 Public Key Infrastructure Time-Stamp Protocol (TSP)}.
    \item \textbf{National Institute of Standards and Technology (NIST) :} \textit{Special Publication 800-63B --- Digital Identity Guidelines: Authentication and Lifecycle Management}.
\end{enumerate}

\end{document}
